%% kapa: exam=1 round=24 year=2013 period=2 form=A subject=appraisal_law subject_ko=부동산관계법규 seq=9 number=49
%% source: Q-Net 감정평가사 기출문제 (cst00309) artlSeq=5073802 file=extracted/1008/5073802/2013년 제24회 감정평가사 1차시험 기출문제/2교시(A형).hwp
%% conversion: kordoc 4.24.51 -> tools/gen1.py (auto); review=proofread-sonnet; flags=none
\begin{question}{49}
국토의 계획 및 이용에 관한 법령상 결정 또는 지정이 실효되는 경우에 해당하는 것은?
\begin{choices}
  \item 지형도면을 따로 작성해야 하는 도시ㆍ군관리계획임에도 불구하고 그 결정의 고시일부터 2년이 되는 날까지 지형도면의 고시가 없는 경우
  \item 도시ㆍ군계획시설결정의 고시일부터 10년이 경과될 때까지 그 시설의 설치에 관한 도시ㆍ군계획시설사업이 시행되지 아니한 경우
  \item 지구단위계획구역 지정에 관한 도시ㆍ군관리계획결정의 고시일부터 1년 이내에 그 지구단위계획구역에 관한 지구단위계획이 결정ㆍ고시되지 아니한 경우
  \item 시가화조정구역의 지정에 관한 도시ㆍ군관리계획결정의 고시일부터 3년이 경과한 경우
  \item 토지거래계약에 관한 허가구역이 지정된 이후 3년간 토지거래계약 허가신청이 없는 경우
\end{choices}
\end{question}
\answer{1}
